\newcommand{\code}[1]{%
  \colorbox{gray!10}{\texttt{\detokenize{#1}}}%
}
\chapter{Formalizing Free XOR in Domino}
\label{chp:domino}
This chapter describes the formalization of the \acrshort{ssp} introduced in the previous chapter using Domino. We first introduce the aspects of Domino related to the formalization, including its project structure and mechanisms used to represent cryptographic constructions and their security arguments. We then describe how the Free XOR construction was translated from its pen-and-paper description to Domino. While the rest of the proof could be successfully verified by Domino, one code equivalence game hop could not be verified. Therefore, we discuss the assumptions and conjectures used, and the invariants and randomness mappings required by the equivalence game hops, as well as the limitations encountered during formalization.
\section{Domino}
\label{sec:whatisdomino}
\newacronym{z3}{Z3}{Z3 theorem prover}
Domino is a proof assistant for State-separating proofs (SSPs) (Section \ref{sec:ssp}) that allows the user to validate the correctness of a SSP using simple and familiar syntax. To allow this, it converts the SSP-style Domino language code into corresponding code in SMT-LIB, a language used by programs called SMT solvers that use a variety of mathematical theorems to determine if a given problem is satisfiable with some possible combination of inputs under specific constraints. If the problem is satisfiable, the SMT solver can also return the inputs that make the problem satisfiable. Domino requires a SMT solver to be installed on the system - either \acrfull{z3} \cite{de_moura_z3_2008}, or CVC5, in order to process these SMT-LIB files that are generated. 

The SSP-style Domino syntax was designed to use very similar keywords and constructs as those in SSPs -- packages, games, and theorems. Packages are composed into games, and these games are used to build a theorem. The explicit relation between the games used in the theorem are defined in the form of game hops. Functional equivalence of pairs of games is argued using invariants that relate their states, and mappings between the randomnesses sampled in either game \cite{domino:paper}. These invariants and randomness mappings are written directly in the SMT-LIB language. An equivalence relation that Domino tries to verify as part of a theorem is referred to as a lemma. In order to establish functional equivalence, games are then checked against one of many in-built and custom lemmata depending on the relations needed in order to prove the correctness of the theorem.

\subsection{The Domino binary}
\label{subsec:architecture}
Domino is published as a Rust crate and is available on GitHub. At the time of writing, Domino is under active development, and its API and language features are subject to frequent changes, including breaking changes between versions.

The main Domino program is exposed as a command line binary. We use two main subcommands in our formalization: 
\begin{itemize}
    \item \code{domino latex} is used to generate \LaTeX  code for the packages, games, theorems and graphs that show how these are related to each other. It uses the underlying \acrshort{smt} solver to generate code that draws the reduction graphs when the output \lstinline|.tex| files are compiled. This command also generates pseudocode for every package that is used in the theorem. \acrshort{z3} is the default \acrshort{smt} solver used in this version of Domino \cite{domino:code}. The \code{.tex} files are stored in the \code{\_build\latex} subdirectory.
    \item \code{domino prove} is used to validate each of the game hops that are defined in the theorem file. It also uses the SMT solver to validate the reduction and equivalence game hops. Equivalence game hops require an additional randomness mapping step, and may require a separate invariant file. \code{domino prove} accepts parameters, and the following are used in this project:
    \begin{itemize}
        \item \code{-t} generates transcripts which can be used to review the outputs of the SMT solver.
        \item \colorbox{gray!10}{\texttt{-\kern0pt-proof}} selects a specifc theorem to be proven and is useful if the proof code is split into multiple theorems and the user needs to prove a specific one. It is followed by a string parameter with the name of the theorem.
        \item \colorbox{gray!10}{\texttt{-\kern0pt-proofstep}} selects the proof step to be proven in the selected theorem. It is followed by a integer which is the index of the proof step to be proven. Indexing starts at 0.
        \item \colorbox{gray!10}{\texttt{-\kern0pt-claim}} selects the lemma to be proven in the selected proofstep. It takes a string parameter that is the full name of the lemma to be proven. 
    \end{itemize}
\end{itemize}  

Domino has two main dependencies:
\begin{itemize}
    \item A recent Rust toolchain. Instructions on setting this up can be found at \cite{rust_install}.
    \item A functioning \acrshort{z3} and CVC5 installation and the executables on \code{PATH}. The installation instructions for these can be found at the respective websites of the projects \cite{cvc5_install,z3_install}.
\end{itemize} 
It is important to note that the current version of Domino prefers \acrshort{z3} for running \code{domino latex} since it is faster at generating the dependency graphs. The \code{-s cvc5} option can be used to force it to use CVC5 for graph resolution.  However, for \code{domino prove}, CVC5 is the default \acrshort{smt} solver, and \acrshort{z3} is only supported using an adaptor script which was not tested by us \cite{domino:paper}.

After setting up these dependencies, Domino can be installed using the installation instructions in the Domino repository \cite{domino:code}.

\subsection{Packages}
\label{subsec:arch_packages}
Packages in Domino are syntactically the same as packages in \acrshort{ssp}s. The core of a package are the oracles that it exposes. Each oracle defines what arguments it expects, and its return values. An SSP package can also be composed with other packages, and Domino allows this using an \code{import oracle} syntax that can be used to declare external oracles that can composed with this package. External oracle calls cannot be inlined, however, and must use an explicit \code{invoke <ORACLE-NAME>} syntax that assigns the value to some variable. Listing \ref{lst:domino-pkg-example} shows a sample Domino package with placeholders.
\lstdefinelanguage{Domino}{
    % morecomment=[s]{/*}{*/},
    moredelim=[is][\color{gray}]{@@}{@@}
}
\begin{lstlisting}[
    language=Domino,
    caption={Example Domino package},
    label={lst:domino-pkg-example},
    commentstyle=\small\rmfamily\itshape,
    basicstyle=\small\ttfamily,
    columns=fixed]
    package Reduction {
    params {
        @@/*an integer parameter*/@@
        n: Integer, 
        @@/*a function parameter that maps two @@
        @@bitstrings of length n to a bitstring of length m*/@@
        fn1: fn Bits(n), Bits(n), -> Bits(m),
    }
    @@/*state of the package*/@@
    state {
        x: Bits(n),  
    }
    @@/*external oracles consumed in this package*/@@
    import oracles {
        EXT1(param1: Integer) -> Bits(m),
    }
    @@/*an oracle exposed by this package*/@@
    oracle ORA(f: Bits(*), x: (Bool, Bool)) -> (Bits(m),Bits(m)) 
    {
        temp <- invoke EXT1(n);
        temp2 <- fn1(n);
        return (temp,temp2);
    }}
\end{lstlisting}


\subsection{Compositions}
\label{subsec:arch_compositions}
Compositions, or games, in Domino are syntactically equivalent to games in \acrshort{ssp}s. A game first defines constants that will be passed to package instances. Then, it defines instances of packages that are used in the definition of the game. The compose block is used to define the composition of one or more packages by declaring the oracles of each package instance that are accessible to other package instances in the composition. The compose block also has a special keyword 'adversary' which is used to define the specific package instances that are accessible to an adversary. Listing \ref{lst:domino-game-example} shows a sample Domino game with placeholders.
% \lstdefinelanguage{Domino}{
%     % morecomment=[s]{/*}{*/},
%     moredelim=[is][\color{gray}]{@@}{@@}
% }
% \begin{lstlisting}[
%     language=Domino,
%     caption={Example Domino game},
%     label={lst:domino-game-example},
%     commentstyle=\small\rmfamily\itshape,
%     basicstyle=\small\ttfamily,
%     columns=fixed]
%     composition Game_R_Pkg1 {
%     @@/*an integer const parameter passed @@
%     @@ to instances*/@@
%     const n: Integer;
%     @@/*a const function that is passed to @@
%     @@ package instances*/@@
%     const fn1: fn Bits(n), Bits(n), -> Bits(m);
%     instance R = Reduction{
%     @@    /*instantiate package Reduction@@
%     @@     with the following parameters*/@@
%         params { n: n, fn1: fn1, }
%     }
%     instance Pkg1 = Package1{
%         params { n: n, }
%     }
%     compose {
%     @@    /*Adversary has access to oracle ORA @@
%     @@     of Reduction instance R*/ @@
%         adversary: { ORA: R },
%     @@    /*R has access to oracle EXT1  @@
%     @@     of Package1 instance Pkg1*/@@
%         R: { EXT1: Pkg1 }
%     }}
% \end{lstlisting}
\begin{longlisting}
\begin{minted}{ssp-lang}
composition Game_R_Pkg1 {
    /*an integer const parameter passed 
     to instances*/
    const n: Integer;
    /*a const function that is passed to 
    package instances*/
    const fn1: fn Bits(n), Bits(n), -> Bits(m);
    instance R = Reduction{
       /*instantiate package Reduction
         with the following parameters*/
        params { n: n, fn1: fn1, }
    }
    instance Pkg1 = Package1{
        params { n: n, }
    }
    compose {
        /*Adversary has access to oracle ORA 
         of Reduction instance R*/ 
        adversary: { ORA: R },
        /*R has access to oracle EXT1  
         of Package1 instance Pkg1*/
        R: { EXT1: Pkg1 }
    }}
\end{minted}
\caption{Example Domino composition}
\label{lst:domino-game-example}
\end{longlisting}
\subsection{Theorems}
\label{subsec:theorems}
Theorems define the relations between different games in Domino. They first define constants that are passed to game instances defined right after. Then, the games that are used in the theorem are instantiated with these constants. Following this, the assumptions block is used to define the assumptions under which the theorem is supposed to be proven. Finally, the gamehops block defines the game hops used to relate the games that the theorem is trying to prove functionally equivalent. These game hops can either be reductions, conjectures, or equivalences. Reductions map two games to each other based on some underlying assumption which is defined in the assumption block earlier. Conjectures declare two games to be functionally equivalent without any further reasoning, and Domino accepts their equivalence as-is. Equivalences map games based on lemmata, and are expressed in terms of randomness mappings and invariant relations. Each equivalence is checked for specific lemmata such as \code{same-output} and \code{equal-aborts} (these two together show input-output equivalence). An equivalence proof step may also contains references to other lemmata that need to be verified for that proof step to be valid. These custom definitions for lemmata are written in SMT-LIB. Listing \ref{lst:domino-theorem-example} shows a sample Domino theorem with placeholders, and without custom lemmata.
% \lstdefinelanguage{Domino}{
%     % morecomment=[s]{/*}{*/},
%     moredelim=[is][\color{gray}]{@@}{@@}
% }
% \begin{lstlisting}[
%     language=Domino,
%     caption={Example Domino theorem},
%     label={lst:domino-theorem-example},
%     commentstyle=\small\rmfamily\itshape,
%     basicstyle=\small\ttfamily,
%     columns=fixed]
%     theorem Theorem1 {
%     @@/*an integer const parameter passed @@
%     @@ to game instances*/@@
%     const n: Integer;
%     @@/*a const function that is passed to @@
%     @@ game instances*/@@
%     const fn1: fn Bits(n), Bits(n), -> Bits(m);
%     @@/*Game instances used in the theorem*/@@
%     instance LeftGame = Game_R_Pkg1{
%         params { n: n, fn1: fn1 }
%     }
%     instance RightGame = Game2{
%         params { n:n, fn1: fn1 }
%     }
%     instance BaseLGame = Game3{
%         params { n:n, fn1: fn1 }
%     }
%     instance BaseRGame = Game4{
%         params { n:n, fn1: fn1 }
%     }
%     @@/*... other games...*/@@
%     assumptions {BaseGame: BaseLGame ~ BaseRGame}
%     gamehops {
%         reduction LeftGame RightGame {
%             assumption BaseGame
%             map {
%                 @@/*map packages of LeftGame and BaseLGame*/@@
%             }
%             map {
%                 @@/*map packages of RightGame and BaseRGame*/@@
%             }
%         }
%         equivalence GameN GameM {
%             invariant: [
%                 ./theorem/invariants/invariant.smt2
%             ]
%             ORACLE: {
%                 lemmas {
%                     invariant: [no-abort]
%                     same-output: [no-abort]
%                     equal-aborts: []
%                     @@/*other lemmata*/@@
%                 }
%             }
%         }
%         @@...@@
%     }}
% \end{lstlisting}
\begin{longlisting}
\begin{minted}{ssp-lang}
theorem Theorem1 {
    /*an integer const parameter passed 
     to game instances*/
    const n: Integer;
    /*a const function that is passed to 
     game instances*/
    const fn1: fn Bits(n), Bits(n), -> Bits(m);
    /*Game instances used in the theorem*/
    instance LeftGame = Game_R_Pkg1{
        params { n: n, fn1: fn1 }
    }
    instance RightGame = Game2{
        params { n:n, fn1: fn1 }
    }
    instance BaseLGame = Game3{
        params { n:n, fn1: fn1 }
    }
    instance BaseRGame = Game4{
        params { n:n, fn1: fn1 }
    }
    /*... other games...*/
    assumptions {BaseGame: BaseLGame ~ BaseRGame}
    gamehops {
    reduction LeftGame RightGame {
        assumption BaseGame
        map {
        /*map packages of LeftGame and BaseLGame*/
        }
        map {
        /*map packages of RightGame and BaseRGame*/
        }
    }
    equivalence GameN GameM {
        invariant: [
            ./theorem/invariants/invariant.smt2
        ]
        ORACLE: {
            lemmas {
                invariant: [no-abort]
                same-output: [no-abort]
                equal-aborts: []
                /*other lemmata*/
            }
        }
    }
    /*...*/
}}
\end{minted}
\caption{Example Domino theorem}
\label{lst:domino-theorem-example}
\end{longlisting}

\subsection{Project Structure}
\label{subsec:projectstructure}
A Domino project directory requires the presence of an empty \code{ssp.toml} file to indicate the project root. The project root contains three other mandatory folders - \code{packages}, \code{games}, and \code{theorem}, each of which contain their namesakes. Most Domino subcommands also take a \code{-{}-path} parameter, followed by a directory, to indicate the project root. 

Each of the three folders in a Domino project can contain one or more \code{.ssp} files. Domino checks packages and games for syntactic correctness even if they are not used in a theorem. This is by design, in order to help packages and games  be checked for syntax before setting up a complicated proof.

\section{Encoding the Free XOR SSP in Domino}
\label{sec:freexorindomino}
Before we describe our encoding of the SSP from the previous chapter in Domino, it is important to note that the pen-and-paper SSP was written in pseudocode, which gives us a lot of freedom on how we represent variables and operations, and allows us to get away with not explicitly defining packages. However, since Domino has its own syntax, and is compiled into SMT-LIB, small syntactical changes in the package code can alter the program's behavior and make it difficult to debug. Therefore, we reason a lot more about packages instead of games (or compositions), and clearly differentiate between them in Domino code. This has the added benefit of allowing us to reflect on the nuances of the Domino programming language and discuss it in Chapter \ref{chp:discussion}. We first discuss our representation of the original pen-and-paper SSP games in Domino, and then describe the other games and game-hops that were used, along with the reasons they were needed. 

\subsection{Common parameters}
The the use of the following paramteres remains consistent across many package and game files, and they are also seen in the theorem file, we discuss them before proceeding.
\begin{itemize}
    \item \code{n}: Parameter used to represent the security parameter $\secpar$.
    \item \code{enc}: Parameter used as an abstract function to represent the encryption function $\enc$.
    \item \code{m}: Parameter used to represent the output length of the function $\enc$. 
    \item \code{b}: Parameter used to switch real and ideal behaviors in a few games. 
    \item \code{xor_}: Parameter used as an abstract function used to represent the bitwise XOR operation, since Domino does not provide a real representation of the bitwise XOR operation. We discuss this in more detail in Section \ref{chp:discussion}.
\end{itemize}

\subsection{The FreeXOR, LIN-RK-KDM, and Reduction games}
\label{subsec:freexorandlinkrkkdmgames}
The reduction described in Chapter \ref{chp:penpaperssp} consists of the games $G^{b}_\mathsf{Sim\text{-}FreeXOR}$ and \allowbreak $G^{b}_{\mathsf{LIN-RK-KDM}}$, and the reduction $\rdv$. It also uses a composition $\rdv\to G^{b}_{\mathsf{LIN-RK-KDM}}$. 
We represent the ideal and real instances of $G^{b}_\mathsf{Sim\text{-}FreeXOR}$ as  two different packages in Domino, namely \code{SimFreeXORIdeal}  and \code{SimFreeXORReal}. Both instances of $G^{b}_{\mathsf{LIN-RK-KDM}}$ are represented as one package \code{LINRKKDM} that takes an input parameter bit \code{b} that can be used to switch between the real and ideal behaviors. 
The reduction $\rdv$ is modeled as its own package \code{Reduction}.
Table \ref{tab:freexor-domino-packages} shows this mapping.

\begin{table}[ht]
    \centering
    \caption{Mapping between the games in the pen-and-paper reduction and Domino packages.}
    \label{tab:freexor-domino-packages}
    \begin{tabularx}{\textwidth}{l l X}
        \toprule
        \textbf{\makecell{Pen-and-paper \\game}} & \textbf{\makecell{Domino representation\\(.pkg.ssp)}} & \textbf{Purpose} \\
        \midrule
        $G^{0}_\mathsf{Sim\text{-}FreeXOR}$ & \texttt{SimFreeXORReal} & The real SimFreeXOR game \\
        $G^{1}_\mathsf{Sim\text{-}FreeXOR}$ & \texttt{SimFreeXORIdeal} & The ideal SimFreeXOR game \\
        $G^{b}_{\mathsf{LIN-RK-KDM}}$ & \texttt{LINRKKDM} & Both \acrshort{linrkkdm} games \\
        $\rdv$ & \texttt{Reduction} & The reduction that will be composed with $G^{b}_{\mathsf{LIN-RK-KDM}}$ \\
        \bottomrule
    \end{tabularx}
\end{table}

However, the Domino code for these packages is not a literal translation of the pseudocode from the pen-and-paper SSP packages. 

\subsubsection{LINRKKDM.pkg.ssp}
\label{subsubsec:linrkkdm}
\begin{longlisting}
\begin{minted}{ssp-lang}
oracle ENC(key: Bits(n), msg: Bits(n), z: Bool) -> Bits(m) {
    if (delta == None) {
        tmp <-$ Bits(n) sample-name delta;
        delta <- Some(tmp);
    } 
    r <-$ Bits(n) sample-name r;
    if b {return enc(xor_(key, Unwrap(delta)), 0_n, r); }
    else {
        if z {
            return enc(xor_(key, Unwrap(delta)), xor_(msg,Unwrap(delta)), r); 
        } else {
            return enc(xor_(key, Unwrap(delta)), msg, r);
}}}
\end{minted}
\caption{Oracle ENC in LINRKKDM.pkg.ssp}
\label{lst:freexor-pkg-linrkkdm-ora-ENC}
\end{longlisting}
In addition to the package parameters, package LINRKKDM also contains a state variable (a bitstring of length \code{n})  called \code{delta} that corresponds to $\Delta$ in the SSP representation. However, in the Domino package, \code{delta} is wrapped in \code{Maybe(...)}. This is because \code{Bits(n)} is not nullable, and therefore cannot be checked for an `unassigned' state, as we do in $G^{b}_{\mathsf{LIN-RK-KDM}}$. Domino provides a nullable type \code{Maybe(...)} that allows us to set \code{delta} to a \code{None} value which we use as a representation for an `unassigned' state.

It is also worth mentioning that for all subsequent games that use the state variable \code{delta}, we continue to use the \code{Maybe(Bits(n))} type instead of raw \code{Bits(n)}. This choice was made to reduce the number of type casts when comparing and relating game states in SMT-LIB (we discuss this in more detail when discussing invariants and randomness mappings). However, the tradeoff is that we need to use an \code{Unwrap} operator when trying to access the value of \code{delta} in the Domino code.

The code for the \code{ENC} oracle in the \code{LINRKKDM} package also differs from the pen-and-paper SSP as it contains an extra conditional branch using \code{if-else} on the boolean variable \code{z} (Listing \ref{lst:freexor-pkg-linrkkdm-ora-ENC}). This is because Domino does not allow us to multiply a boolean value $z$ with a bitstring $\Delta$, and we cannot represent ${msg} \xor {z}\cdot \Delta$ from $G^{b}_{\mathsf{LIN-RK-KDM}}$ directly.

A final observation is that the \code{enc} parameter that represents the $\enc$ operation from $G^{b}_{\mathsf{LIN-RK-KDM}}$ contains three parameters instead of two. Because $\enc$ is a non-deterministic function, and the abstract functions we define in Domino are treated as deterministic, we introduce a third parameter that takes a freshly sampled random value to emulate the desired random behaviour.

\subsubsection{Reduction.pkg.ssp}
\label{subsubsec:reduction}
%\lstdefinelanguage{Domino}{
    % morecomment=[s]{/*}{*/},
%     moredelim=[is][\color{gray}]{@@}{@@}
% }
% \begin{lstlisting}[
%     language=Domino,
%     caption={Part of oracle GARBLEAND in Reduction.pkg.ssp},
%     label={lst:freexor-pkg-reduction-ora-GARBLEAND},
%     commentstyle=\small\rmfamily\itshape,
%     basicstyle=\small\ttfamily,
%     columns=fixed,
%     frame=single,
%     breaklines=true,
%     float,floatplacement=H]
% oracle GARBLEAND(x: (Bool, Bool)) -> ( Table((Bits(m),Bits(m)), ()), (Bits(n), Bits(n)), Table(Bits(n), Bool)) {
%     @@...@@
%     if ((a==false) and (b==false)) {
%         c <- false; /* AND gate */
%         A_0 <-tempsample1;
%         B_0 <-tempsample2;
%         C_0 <-tempsample3;
%         GarbledRow0 <- (enc(A_0, R_00, r01), enc(B_0, xor_(R_00, C_0), r02));
%         temp <- invoke ENC(B_0, xor_(R_01, C_0), false);
%         GarbledRow1 <- (enc(A_0, R_01, r11), temp);
%         temp <- invoke ENC(A_0, xor_(R_10, C_0), false);
%         GarbledRow2 <- (temp, enc(B_0, R_10, r22)); 
%         temp <- invoke ENC(A_0, R_11, false);
%         temp2 <- invoke ENC(B_0, xor_(R_11, C_0), true);
%         GarbledRow3 <- (temp, temp2);
%         C_[GarbledRow0] <- Some(());
%         C_[GarbledRow1] <- Some(());
%         C_[GarbledRow2] <- Some(());
%         C_[GarbledRow3] <- Some(());
%         X_ <- (A_0, B_0);
%         d[C_0] <- Some(c); 
%     }
%     @@...@@
%     return (C_,X_,d);
% } 
% \end{lstlisting}
\begin{longlisting}
\begin{minted}{ssp-lang}
oracle GARBLEAND(x: (Bool, Bool)) -> 
(Table((Bits(m),Bits(m)), ()), (Bits(n), Bits(n)), Table(Bits(n), Bool)) {
    /*...*/
    if ((a==false) and (b==false)) {
        c <- false; /* AND gate */
        A_0 <-tempsample1;
        B_0 <-tempsample2;
        C_0 <-tempsample3;
        GarbledRow0 <- (enc(A_0, R_00, r01), enc(B_0, xor_(R_00, C_0), r02));
        temp <- invoke ENC(B_0, xor_(R_01, C_0), false);
        GarbledRow1 <- (enc(A_0, R_01, r11), temp);
        temp <- invoke ENC(A_0, xor_(R_10, C_0), false);
        GarbledRow2 <- (temp, enc(B_0, R_10, r22)); 
        temp <- invoke ENC(A_0, R_11, false);
        temp2 <- invoke ENC(B_0, xor_(R_11, C_0), true);
        GarbledRow3 <- (temp, temp2);
        C_[GarbledRow0] <- Some(());
        C_[GarbledRow1] <- Some(());
        C_[GarbledRow2] <- Some(());
        C_[GarbledRow3] <- Some(());
        X_ <- (A_0, B_0);
        d[C_0] <- Some(c); 
    }
    /*...*/
    return (C_,X_,d);
} 
\end{minted}
\caption{Part of oracle GARBLEAND in Reduction.pkg.ssp}
\label{lst:freexor-pkg-reduction-ora-GARBLEAND}
\end{longlisting}
The \code{Reduction} package begins with the package parameters and does not use a state variable, which is consistent with its pen-and-paper counterpart. However, since $\rdv$ is composed with $G^{b}_{\mathsf{LIN-RK-KDM}}$, we use the \code{import oracle} syntax in \code{Reduction} to declare that this package is intended to be composed with another package that exposes a \code{ENC} oracle with a specific signature. This also allows us to invoke the \code{ENC} oracle in the oracles of \code{Reduction}. 

\code{Reduction} exposes one oracle -- \code{GARBLEAND}, which maps to the $\mathsf{GARBLE}$ oracle of $\rdv$, but does not take $f$ as a paramter. This is because Domino does not allow equivalence relations between games that expose oracles with different oracle signatures, and one of the latter games in the equivalence relations we try to establish (specifically $G^{1}_\mathsf{Sim\text{-}FreeXOR}$ and its corresponding package \code{SimFreeXORIdeal}) require a modified oracle signature. We discuss this in more detail later in this subsection. We are able to get around not having $f$ as an input parameter since we already know that we are dealing with the garbling of an AND gate.

Furthermore, since Domino support neither assignment operations to arbitrary boolean indexes (such as $A_a \sample \bin^\secpar $), nor operations between boolean and bitstrings, we got around it using conditional branching on the boolean variables \code{a} and \code{b}. This enabled us to compute the value of \code{c} or $f(x)$ in each branch separately by setting the value of \code{c} to a precomputed $a\land b$.

talk about
\begin{itemize}
    \item enc random variables
    \item sampling outside the conditional branches to simplify randomness mapping later
    \item Some(()) for table value - used to emulate a set instead of an actual table
\end{itemize}

% \lstdefinelanguage{Domino}{
%     % morecomment=[s]{/*}{*/},
%     moredelim=[is][\color{gray}]{@@}{@@}
% }
% \begin{lstlisting}[
%     language=Domino,
%     caption={Importing Oracle \code{ENC}},
%     label={lst:importoracle},
%     commentstyle=\small\rmfamily\itshape,
%     basicstyle=\small\ttfamily,
%     columns=fixed]
%     package Reduction {
%     @@ ... @@
%     import oracles {
%         ENC(key: Bits(n), msg: Bits(n), z: Bool) -> Bits(m),}
%     @@ ... @@
%     }
% \end{lstlisting}
\begin{longlisting}
\begin{minted}{ssp-lang}
package Reduction {
    /*...*/
    import oracles {
        ENC(key: Bits(n), msg: Bits(n), z: Bool) -> Bits(m),}
    /*...*/
    }
\end{minted}
\caption{Importing Oracle \code{ENC} in \code{Reduction}}
\label{lst:importoracle}
\end{longlisting}
We also define the \code{GARBLE} oracle which encodes the code from $\mathsf{GARBLE}$ in the SSP. While randomly sampling variables, we can assign them a sample name, which is useful when we want to relate these randomnesses to the randomnesses generated in other packages and games. 

\subsubsection{SimFreeXORIdealANDONLY.pkg.ssp}
We were able to encode $G^1_\mathsf{Sim\text{-}FreeXOR}$ game without making any significant changes.

\subsubsection{SimFreeXORRealANDONLY.pkg.ssp}
We were able to encode $G^0_\mathsf{Sim\text{-}FreeXOR}$ game using only the modifications we described in the above package files.
\subsection{Games}
\label{subsec:encode_games}
Now, we move on to the mapping of the games in the SSP to the games that we used in Domino. Table \ref{tab:freexor-domino-games} shows this mapping.
\begin{table}[ht]
    \centering
    \caption{Mapping between the games in the pen-and-paper reduction and their Domino representations.}
    \label{tab:freexor-domino-games}
    \begin{tabularx}{\textwidth}{l l X}
        \toprule
        \textbf{\makecell{Pen-and-paper \\game}} & \textbf{\makecell{Domino representation\\(.comp.ssp)}} & \textbf{Purpose} \\
        \midrule
        $G^{0}_\mathsf{Sim\text{-}FreeXOR}$ & \texttt{GSimFreeXORReal} & The real SimFreeXOR game \\
        $G^{1}_\mathsf{Sim\text{-}FreeXOR}$ & \texttt{GSimFreeXORIdeal} & The ideal SimFreeXOR game \\
        $G^{b}_{\mathsf{LIN-RK-KDM}}$ & \texttt{GLINRKKDM} & Both \acrshort{linrkkdm} games \\
        $\rdv\to G^{b}_{\mathsf{LIN-RK-KDM}}$ & \texttt{RcompGLINRKKDM.comp} & The reduction composed with $G^{b}_{\mathsf{LIN-RK-KDM}}$ \\
        \bottomrule
    \end{tabularx}
\end{table}

\subsubsection{GLINRKKDM.comp.ssp}
\code{GLINRKKDM.comp.ssp} encodes the LIN-RK-KDM game. It first defines the constants that the package it instantiates needs. Then, it creates an instance of the \code{LINRKKDM} package and passes those parameters to it. Finally, it defines the oracle interface that it exposes to an adversary.

\subsubsection{RcompGLINRKKDM.comp.ssp}
\code{RcompGLINRKKDM.comp.ssp} encodes the composition of the reduction $\rdv$ with the LIN-RK-KDM game. While following the same syntax as the previous game file, it instantiates two different games, and composes them by defining that the \code{ENC} oracle of the \code{LINRKKDM} instance is available to the instance of the \code{Reduction} package. It also defines the interface exposed to the adversary.

\subsubsection{GSimFreeXORReal.comp.ssp and GSimFreeXORIdeal.comp.ssp}
These two files follow the same syntax as the previous game files and define their respective games.

\subsection{Theorems}
\label{subsec:encode_therorems}
Theorem files are where we define the relationships between the games we have defined so far.
We start by defining the constant variables that we would like to instantiate the games with. These parameters are then injected into the games (and then their packages), making their way to the underlying oracles. This is a useful property to have if we want some functions or variables to behave the same way in every package.
This allows us to use the same parameters for each package, causing each package to have the same message length, security parameter, and for them to use the same encryption scheme and XOR functions.

Then, we instantiate all the games associated with this theorem, and separate the real and ideal games by passing relevant parameters, in the case of games that have varying logic based on an input paramter bit, like in the case of game \code{GLINRKKDM}. 
After creating instances of all the games, we define the assumptions that we would like to reduce our games to in future reduction game hops. In our case, this is Definition \ref{def:linrkkdmsecurity} (LIN-RK-KDM security). We define our assumption as the equivalence of the instance of game \code{GLINRKKDM} initialized with bit 0, and its instance initialized with bit 1.

Finally, we define the game hops used in the theorem. Game hops can be of two types, reductions and equivalences.
\subsubsection{Reductions}
A reduction game hop is defined by first choosing the games whose security has to be reduced to an assumption. Following Theorem \ref{thm:reduction}, we choose to reduce the instances of game \code{RcompGLINRKKDM} instantiated with bit 0 and 1 to the \code{LINRKKDM} game. After choosing the composed games that need to be reduced to the  assumption games, we must then manually map the package instances that both composed game have to their corresponding packages in the assumption games respectively. Listing \ref{lst:domino-theorem-reduction-gamehop} shows the reduction game hop we use in this theorem.
% \lstdefinelanguage{Domino}{
%     % morecomment=[s]{/*}{*/},
%     moredelim=[is][\color{gray}]{@@}{@@}
% }
% \begin{lstlisting}[
%     language=Domino,
%     caption={Reduction game hop in the theorem},
%     label={lst:domino-theorem-reduction-gamehop},
%     commentstyle=\small\rmfamily\itshape,
%     basicstyle=\small\ttfamily,
%     columns=fixed]
%     theorem freexor {
%     @@ ... @@
%     assumptions {
%         lin_rk_kdm: GLINRKKDM0 ~ GLINRKKDM1
%     }
%     gamehops{
%     reduction RcompGLINRKKDM0 RcompGLINRKKDM1 {
%             assumption lin_rk_kdm
%             map GLINRKKDM0 RcompGLINRKKDM0 {
%                 GLINRKKDM: GLINRKKDM
%             }
%             map GLINRKKDM1 RcompGLINRKKDM1 {
%                 GLINRKKDM: GLINRKKDM
%             }
%     }}}
% \end{lstlisting}
\begin{longlisting}
\begin{minted}{ssp-lang}
theorem freexor {
    /*...*/
    assumptions {
        lin_rk_kdm: GLINRKKDM0 ~ GLINRKKDM1
    }
    gamehops{
    reduction RcompGLINRKKDM0 RcompGLINRKKDM1 {
        assumption lin_rk_kdm
        map GLINRKKDM0 RcompGLINRKKDM0 {
            GLINRKKDM: GLINRKKDM
        }
        map GLINRKKDM1 RcompGLINRKKDM1 {
            GLINRKKDM: GLINRKKDM
        }
}}}
\end{minted}
\caption{Reduction game hop in the theorem}
\label{lst:domino-theorem-reduction-gamehop}
\end{longlisting}

\subsubsection{Equivalences}
%Since invariants and randomness mappings are part of equivalences, they go here 
A code equivalence argument can be encoded using an equivalence game hop in Domino. First, we choose the oracles exposed by the two games (to the adversary) for which we would like to show equivalence. Then, we describe what equivalence actually means in this context using varying lemmata. Domino offers two built-in lemmata - \code{same-output} and \code{equal-aborts}. To put it simply, the \code{same-output} lemma, along with the \code{[no-abort]} argument verifies whether the oracles of the two games return the same output when they are provided the same inputs. The \code{equal-abort} lemma verifies if both oracles also abort on the same inputs, verifying equivalence of their failure behavior too. 

We also have extra conditions that are required for our proof to be validated, which can be defined as extra lemmata for Domino to verify. We define these in the file \code{real\_game\_invariant.smt2}, and is discussed in the following section.
Listing \ref{lst:domino-theorem-equivalence-gamehop} shows one of the equivalence game hops we use in this theorem.
% \lstdefinelanguage{Domino}{
%     % morecomment=[s]{/*}{*/},
%     moredelim=[is][\color{gray}]{@@}{@@}
% }
% \begin{lstlisting}[
%     language=Domino,
%     caption={Equivalence game hop in the theorem},
%     label={lst:domino-theorem-equivalence-gamehop},
%     commentstyle=\small\rmfamily\itshape,
%     basicstyle=\small\ttfamily,
%     columns=fixed]
%     theorem freexor {
%     @@ ... @@
%     gamehops{
%     equivalence RcompGLINRKKDM0 GSimFreeXORReal {
%             GARBLE: {
%                 invariant: [
%                     ./theorem/real_game_invariant.smt2
%                 ]
%                 lemmas {
%                     invariant: [no-abort]
%                     equal-aborts: []
%                     same-output: [no-abort]
%                 }
%     }}}}
% \end{lstlisting}
\begin{longlisting}
\begin{minted}{ssp-lang}
theorem freexor {
    /*...*/
    gamehops{
    equivalence RcompGLINRKKDM0 GSimFreeXORReal {
            GARBLE: {
                invariant: [
                    ./theorem/real_game_invariant.smt2
                ]
                lemmas {
                    invariant: [no-abort]
                    equal-aborts: []
                    same-output: [no-abort]
                }
    }}}}
\end{minted}
\caption{Equivalence game hop in the theorem}
\label{lst:domino-theorem-equivalence-gamehop}
\end{longlisting}

\subsection{Invariants and Randomness Mappings}
Domino allows users to write custom lemmata that can be verified when \code{domino} \code{prove} runs, making it very extensible. We can declare such lemmata by attaching their corresponding \code{.smt2} file, and adding it to the list of lemmata to be proven. We do this in Listing \ref{lst:domino-theorem-equivalence-gamehop} by declaring a lemma called \code{invariant}. 
These extension files are written in SMT-LIB. Listing \ref{lst:domino-theorem-equivalence-invariant}, is an excerpt from the invariant file we use. The first function that the file defines is a randomness mapping between the left and right games. Randomness mappings are necessary because it allows us to ensure that the randomness sampled in both games are related. This is important when our oracle outputs depend on the random values, and equivalence conditions will simply fail unless the randomnesses are defined to be the same.
The other function is the invariant, and we define a condition here that we need to be true for the equivalence to be established. Specifically, the circuit offset $\Delta$ chosen in both games needs to be the same, otherwise the oracles will return different outputs and the \code{same-output} property will never be satisfied.
% \begin{lstlisting}[
%     language=Lisp,
%     caption={Invariant and randomness mapping used in the theorem},
%     label={lst:domino-theorem-equivalence-invariant},
%     commentstyle=\small\rmfamily\itshape,
%     basicstyle=\small\ttfamily,
%     columns=fixed]
%     (define-fun randomness-mapping-GARBLE
%     (
%         (sample-id-left SampleId)
%         (sample-id-right SampleId)
%         (sample-offset-left Int)
%         (sample-offset-right Int)
%     )
%     Bool
%     (or
%         (and
%             (= sample-id-left (sample-id "R" "GARBLE" "R_00"))
%             (= sample-id-right (
%                 sample-id "SimFreeXOR0" 
%                 "GARBLE" "R_00")
%             )
%             (= sample-offset-left 0)
%             (= sample-offset-right 0)
%         )
%         (and
%             (= sample-id-left (
%                 sample-id "LINRKKDM" 
%                 "ENC" "delta")
%             )
%             (= sample-id-right (
%                 sample-id "SimFreeXOR0" 
%                 "GARBLE" "delta")
%             )
%             (= sample-offset-left 0)
%             (= sample-offset-right 0)
%         )
%     )
%     )
%     (define-state-relation invariant
%     (left-game right-game) 
%     (and
%         (= 
%             (maybe-get left-game.GLINRKKDM.delta) 
%             (right-game.SimFreeXOR0.delta))
%     )
%     )
% \end{lstlisting}
\begin{longlisting}
\begin{minted}{smt2-lang}
(define-fun randomness-mapping-GARBLE
    (
        (sample-id-left SampleId)
        (sample-id-right SampleId)
        (sample-offset-left Int)
        (sample-offset-right Int)
    )
    Bool
    (or
        (and
            (= sample-id-left (sample-id "R" "GARBLE" "R_00"))
            (= sample-id-right (
                sample-id "SimFreeXOR0" 
                "GARBLE" "R_00")
            )
            (= sample-offset-left 0)
            (= sample-offset-right 0)
        )
        (and
            (= sample-id-left (
                sample-id "LINRKKDM" 
                "ENC" "delta")
            )
            (= sample-id-right (
                sample-id "SimFreeXOR0" 
                "GARBLE" "delta")
            )
            (= sample-offset-left 0)
            (= sample-offset-right 0)
        )
    )
    )
    (define-state-relation invariant
    (left-game right-game) 
    (and
        (= 
            (maybe-get left-game.GLINRKKDM.delta) 
            (right-game.SimFreeXOR0.delta))
    )
)
\end{minted}
\caption{Invariant and randomness mapping used in the theorem}
\label{lst:domino-theorem-equivalence-invariant}
\end{longlisting}
\subsection{Verification Results}
\label{subsec:results}
The resulting Domino formalization was tested by using the verification mechanism provided in Domino (\code{domino prove}). The reduction between \code{RcompGLINRKKDM0} and \code{RcompGLINRKKDM1} was successfully verified under the LIN-RK-KDM assumption. However, the equivalence relation between \code{RcompGLINRKKDM0} and \code{GSimFree-} \code{XORReal}, and the equivalence relation between \code{RcompGLINRKKDM1} and \code{GSimFree-} \code{XORIdeal} could not be successfully established. This suggests the requirement of more invariants, which we discuss in Chapter \ref{chp:discussion}.

The full code for this formalization can be found on GitHub \cite{prateekfreexor:code}.
